\documentclass[11pt]{article}
\usepackage{mathblatt}
% Präambel des Setzers (werkzeuge/setzer-praeambel.tex, wortgleich); dazu \gteil (Muster T2U)
\makeatletter
\setlength{\mbpfbe}{0mm}% keine Punkte (pruefung.md § 5: nur beim Prüfstein)
\newlength{\gEin}\setlength{\gEin}{0pt}
\let\gSchrift\relax
\renewcommand{\pfkreuzzeile}[1]{\par\noindent\hangindent3.2em\hangafter1\hspace*{1.6em}$\square$\ #1\par}
\newcommand{\gkreuz}[1]{$\square$\ #1\hspace{2em}}
\newcommand{\gkreuzzeile}[2]{\par\noindent\hangindent2.2em\hangafter1\mbox{#1\,$\square$}\ #2\par}% Buchstabe, Option: eine je Zeile (neben Skizze)
% links, was man ansieht, rechts, was man tut; Spaltengrenze überall an derselben Stelle (gemeinsam.md „Satz“)
% v0.17: Fußeintrag „Nr: Wert“ – kurz in einer Box (nie getrennt), lang umbrechbar mit der Nummer am Anfang
% (vorher \mbox: lange Lösungen liefen über den Rand, AS40 138 pt); robust, weil \fusshilfe in die .aux schreibt
\NewDocumentCommand{\gfuss}{m m}{\sbox0{#1:~#2}\ifdim\wd0<0.45\textwidth\mbox{#1:~#2}\else\mbox{#1:}\nobreak\ #2\fi}
\newsavebox{\gbox}\newlength{\gLinks}
% v0.10 (Lehrer 13.10.): keine starre Halbierung – die linke Spalte ist so breit wie die Figur (+ 6 mm), die
% Antwort steht direkt daneben; ist die Figur breiter als 0,62 Zeile, steht die Antwort darunter
\newcommand{\gzwei}[2]{\par\smallskip\noindent\sbox{\gbox}{#1}%
  \setlength{\gLinks}{\dimexpr\wd\gbox+6mm\relax}%
  \ifdim\gLinks>0.62\linewidth
    \usebox{\gbox}\par\smallskip\noindent\begin{minipage}{\linewidth}#2\end{minipage}%
  \else
    \begin{minipage}[t]{\gLinks}\vspace{0pt}\usebox{\gbox}\end{minipage}%
    \begin{minipage}[t]{\dimexpr\linewidth-\gLinks\relax}\vspace{0pt}\raggedright #2\end{minipage}%
  \fi\par}
% v0.9 (12.10., pruefung.md § 5 „Marke mit B, Bedingungen rechts“): drittes Argument von pfaufg = Bedingungen
% der Prüfung (v0.10: \pfohnetr = TR durchgestrichen, „FOR“), klein grau im rechten Rand auf der Grundlinie der Nummer, Zeilen mit \\
\newcommand{\gbedrand}[2]{\ifblank{#2}{}{\rlap{\hspace*{\dimexpr#1+2mm\relax}%
  \parbox[t]{15mm}{\raggedright\scriptsize\leavevmode\color{mbgrau}\strut#2}}}}
% Teil eines Bündels (a, b …): Marke links im Rand, Bedingungen rechts im Rand, auf der Zeile des Teils
\newcommand{\gteilrand}[2]{\llap{\mbpfrandmarke{#1}\hspace*{\mbpfnr}}\gbedrand{\linewidth}{#2}}
% Teil eines Bündels: Buchstabe in eigener Spalte, Inhalt eingerückt (\gEin = Einzug, damit die Spaltengrenze von
% \gzwei an derselben Stelle bleibt)
\newenvironment{gbteil}[3]{\par\smallskip\noindent\gteilrand{#2}{#3}\makebox[\mbpfteil][l]{\textbf{#1}}%
  \begin{minipage}[t]{\dimexpr\linewidth-\mbpfteil\relax}\setlength{\gEin}{\mbpfteil}\raggedright\strut\ignorespaces}%
 {\end{minipage}\par}
% Grafik höchstens so breit wie die Spalte (Bündel nebeneinander, Koordinatensystem ohne Karo)
\newsavebox{\gpbox}
\newcommand{\gpassend}[1]{\sbox{\gpbox}{#1}\ifdim\wd\gpbox>\linewidth\typeout{GPASSEND: verkleinert \the\wd\gpbox\space auf \the\linewidth}\resizebox{\linewidth}{!}{\usebox{\gpbox}}\else\usebox{\gpbox}\fi}
% Skizzenfeld (v0.10, Lehrer 13.10.): Karo ab dem linken Rand über die Zeile, „Skizze“ klein grau oben links im Feld
% (wie \pfskizzenfeld der Vorlage); der rechte Rand bleibt den Bedingungen
\newcommand{\gskizze}[1]{\par\smallskip\noindent\begin{tikzpicture}
  \clip (0,0) rectangle ({\linewidth-1mm},{#1*10mm});
  \draw[black!16,line width=0.3pt,step=5mm] (0,0) grid ({\linewidth},{#1*10mm});
  \node[anchor=north west,inner sep=2pt,font=\scriptsize,text=mbgrau,fill=white] at (0,{#1*10mm}) {Skizze};
  \end{tikzpicture}\par}
\RenewDocumentEnvironment{pfaufg}{m m m}{%
  \begin{lrbox}{\mbaufgabebox}\begin{minipage}[t]{\linewidth}%
  \gSchrift\noindent\hspace*{\gEin}\mbpfrandmarke{#1}\gbedrand{\linewidth-\gEin}{#3}%
  \begin{minipage}[t]{\mbpfnr}\strut\textbf{#2}\end{minipage}%
  \begin{minipage}[t]{\dimexpr\linewidth-\gEin-\mbpfnr\relax}\raggedright\strut\ignorespaces}%
 {\end{minipage}%
  \end{minipage}\end{lrbox}%
  \setlength{\mbaufgabehoehe}{\dimexpr\ht\mbaufgabebox+\dp\mbaufgabebox\relax}%
  \ifdefempty{\gkopfnach}{\par\addvspace{5pt}\Needspace*{\mbaufgabehoehe}}{%
    % Stufenkopf, der mit dieser (ersten) Aufgabe zusammenbleibt: Platz für beide verlangen, dann den Kopf setzen
    % (12.10., Quadratische: ein Kopf stand allein unten, die Aufgabe mit Koordinatensystem auf der nächsten Seite)
    % v0.17 (AS40 „Terme deuten“ allein unten): hinter dem Kopf steht das Sprungziel (\special, ein Knoten, nach
    % dem die Seite brechen durfte) – darum \nobreak nach dem Kopf; Vorhalt 56 pt = Kopf mit Linie und Abständen
    \par\Needspace*{\dimexpr\mbaufgabehoehe+56pt\relax}\gkopfnach\par\nobreak\gdef\gkopfnach{}\addvspace{5pt}}%
  \noindent\usebox{\mbaufgabebox}\par\addvspace{5pt}}
% Stufenkopf vormerken: \gkopfvor{<Kopf>} – die nächste pfaufg setzt ihn, sobald beide zusammen Platz haben
\def\gkopfnach{}
\newcommand{\gkopfvor}[1]{\gdef\gkopfnach{#1}}
% Teilaufgabe mit Skizze links in gzwei: Buchstabe in eigener Spalte, darunter Text und Skizze
\newcommand{\gteil}[2]{\makebox[\mbpfteil][l]{\textbf{#1}}\parbox[t]{\dimexpr0.5\textwidth-\gEin-\mbpfnr-5mm-\mbpfteil\relax}{\raggedright #2}}
% Stufenkopf der Vorlage (mathblatt.sty 2026-10-12: Linie, Abstand, 13 pt, \pfjahre) ohne ihren Vorhalt von
% 14 Zeilen: der Setzer fordert vor dem Kopf 0,16 Seitenhöhe (Kopf + erste Aufgabe), mehr nicht
\setlength{\mbpfstufevorhalt}{0pt}
\makeatother
\newcommand{\karo}[1]{\par\smallskip\noindent\begin{tikzpicture}
  \clip (0,0) rectangle ({\linewidth-1mm},{#1*10mm});
  \draw[black!16,line width=0.3pt,step=5mm] (0,0) grid ({\linewidth},{#1*10mm});
  \end{tikzpicture}\par}
\newcommand{\kk}[1]{$\square$\,#1\hspace{0.9em}}
\newcommand{\linie}[1]{\,{\color{black!45}\rule[-1pt]{#1}{0.4pt}}\,}
\newcommand{\pkt}{\textperiodcentered{}}
% v0.14 (14.10.): Stufenkopf mit fünf Punkten (Option = Jahre für die Bestellkarte) oder „zuletzt 19“
\NewDocumentCommand{\gpunkte}{o m}{#2}
\newcommand{\gzuletzt}[1]{zuletzt #1}
% v0.7 (12.10.): Schülerfrage unter dem Stufenkopf; Ankreuzoption als Bild (Buchstabe und Kästchen darunter)
\newcommand{\gfrage}[1]{\par\nobreak\noindent\hspace*{\mbpfnr}{\small\itshape #1}\par\nobreak\addvspace{2pt}}
\newlength{\gbildbreite}\setlength{\gbildbreite}{0.32\linewidth}% Setzer setzt sie je Reihe (vier Bilder: 0,24)
\newcommand{\gbildoption}[3]{\begin{minipage}[t][#3][b]{\gbildbreite}\centering #2\par\smallskip\mbox{#1\,$\square$}\end{minipage}}
% Lineare (12.10.): Bild mit Buchstabe ohne Kästchen (Diagramm zuordnen); Aussage mit „wahr / falsch“ zum Ankreuzen
\newcommand{\gbildname}[3]{\begin{minipage}[t][#3][b]{\gbildbreite}\centering #2\par\smallskip\mbox{\textbf{#1}}\end{minipage}}
\newcommand{\gwahrfalsch}[2]{\par\noindent\hangindent2.2em\hangafter1\makebox[2.2em][l]{(#1)}#2\hfill\mbox{wahr $\square$\quad falsch $\square$}\par}
\begin{document}

\pfheftstil
\pflaufkopf[Abitur GK \pkt{} Prüfung]{Winkel am Graphen}{AA80}\pfseitevon
\pfheftkopf{Winkel am Graphen}{}
\pfstufenzeile{Steigungswinkel gesucht \pkt{} Schnittwinkel gesucht}

% ===================== Stufe: Steigungswinkel gesucht =====================
\gkopfvor{\pfstufekopf{Steigungswinkel gesucht}{\gpunkte[25 24 22]{\textcolor{black!25}{$\circ$}\,\textcolor{black!80}{$\bullet$}\,\textcolor{black!80}{$\bullet$}\,\textcolor{black!25}{$\circ$}\,\textcolor{black!80}{$\bullet$}}}\special{pdf:dest (stufe-steigungswinkel-gesucht) [@thispage /XYZ @xpos @ypos null]}}
% kein Einstieg: keine Schülerfrage in der Gliederung (Schülerfrage, -Fälle, -Form, -Weg, -Skizze; oder Frage) – Einstieg nur von dort, nie aus der Bank (pruefung.md § 5, 12.10.)
% Teile eines Originals (2022-bebb-gk-B2.2): 2022-bebb-gk-B2.2e-m1, 2022-bebb-gk-B2.2f-m1
\begin{pfaufg}{}{1.}{}
% 2022-bebb-gk-B2.2e-m1 (Original 2022-bebb-gk-B2.2e, Schwere II, rang 23; herausgelöste Fassung: keine)
\vspace*{-\baselineskip}\vspace*{-\smallskipamount}
\begin{gbteil}{a)}{\raisebox{\dimexpr-\height+\ht\strutbox\relax}{\shortstack[r]{GK~’22\\B2.2e}}}{}
Gegeben ist die Funktion $f$ mit $f(x) = -\tfrac{1}{18}x^3 + \tfrac{1}{3}x^2$, $x \in \mathbb{R}$. Ihr Graph hat den Wendepunkt $W\left(2 \mid \tfrac{8}{9}\right)$. Ermittle den Winkel, unter dem die Tangente in $W$ die $x$-Achse schneidet. Gib an, bei welchem $y$-Wert diese Tangente die $y$-Achse schneidet.
\karo{4}
\end{gbteil}
% 2022-bebb-gk-B2.2f-m1 (Original 2022-bebb-gk-B2.2f, Schwere II, rang 24; herausgelöste Fassung: keine)
\begin{gbteil}{b)}{\raisebox{\dimexpr-\height+\ht\strutbox\relax}{\shortstack[r]{GK~’22\\B2.2f}}}{}
Für die Funktion $f$ mit $f(x) = -\tfrac{1}{18}x^3 + \tfrac{1}{3}x^2$ gilt: Die Tangente an den Graphen an der Stelle $x_T$ steigt genau dann unter einem Winkel von mindestens 31,0°, wenn $a \le x_T \le b$ ist. Ermittle $a$ und $b$ (auf zwei Nachkommastellen gerundet).
\karo{4}
\end{gbteil}
\fusshilfe{\gfuss{1}{a)~33,7°; y-Achse bei $-\tfrac{4}{9}$; b)~2,63}}
\end{pfaufg}
% 2024-bebb-gk-B2.1i-m1 (Original 2024-bebb-gk-B2.1i, Schwere III, rang 25; herausgelöste Fassung: keine)
\begin{pfaufg}{GK~’24\\B2.1i}{2.}{}
Ein Schrankknauf entsteht durch Rotation des Graphen von $f$ mit $f(x) = (x - 3) \cdot e^{-\tfrac{x}{3} + 2}$, $3 \le x \le 12$, um die $x$-Achse; es gilt $f'(x) = \left(2 - \tfrac{x}{3}\right) \cdot e^{-\tfrac{x}{3} + 2}$. Im Querschnitt bilden die Graphen von $f$ und $-f$ an der Spitze $S(3 \mid 0)$ den Winkel $\alpha$. Vorgabe: $\alpha$ soll mindestens 135° betragen. Prüfe rechnerisch, ob die Vorgabe eingehalten wird.
\karo{4}
\fusshilfe{\gfuss{2}{Vorgabe eingehalten}}
\end{pfaufg}
% 2025-bebb-gk-B2.1d-m1 (Original 2025-bebb-gk-B2.1d, Schwere III, rang 26; herausgelöste Fassung: keine)
\begin{pfaufg}{GK~’25\\B2.1d}{3.}{}
Gegeben ist die Funktion $f$ mit $f(x) = \tfrac{1}{18}x^3 - \tfrac{3}{2}x$, $x \in \mathbb{R}$, und die Gerade $g$ mit $y = \tfrac{1}{2}x$. Die Tangente $k$ an den Graphen von $f$ im Punkt $\left(3\sqrt{3} \mid 0\right)$ schneidet $g$ unter dem Winkel $\alpha$. Bestimme den Steigungswinkel von $k$. Untersuche, ob $\sin \alpha = \cos \alpha$ gilt.
\karo{4}
\fusshilfe{\gfuss{3}{71,6°; cos $\alpha$}}
\end{pfaufg}

% ===================== Stufe: Schnittwinkel gesucht =====================
\gkopfvor{\pfstufekopf{Schnittwinkel gesucht}{\gpunkte[23 22]{\textcolor{black!25}{$\circ$}\,\textcolor{black!25}{$\circ$}\,\textcolor{black!25}{$\circ$}\,\textcolor{black!80}{$\bullet$}\,\textcolor{black!80}{$\bullet$}}}\special{pdf:dest (stufe-schnittwinkel-gesucht) [@thispage /XYZ @xpos @ypos null]}}
% kein Einstieg: keine Schülerfrage in der Gliederung (Schülerfrage, -Fälle, -Form, -Weg, -Skizze; oder Frage) – Einstieg nur von dort, nie aus der Bank (pruefung.md § 5, 12.10.)
% 2022-bebb-gk-B2.1e-m1 (Original 2022-bebb-gk-B2.1e, Schwere II, rang 27; herausgelöste Fassung: keine)
\begin{pfaufg}{GK~’22\\B2.1e}{4.}{}
Gegeben ist die in $\mathbb{R}$ definierte Funktion $f$ mit $f(x) = (x + 3) \cdot e^{-x}$ und ihrer Ableitung $f'(x) = -(x + 2) \cdot e^{-x}$. Die Graphen von $f$ und $f'$ schneiden sich an der Stelle $-2{,}5$. Ermittle den Winkel, unter dem sie sich schneiden.
\karo{4}
\fusshilfe{\gfuss{4}{12,5°}}
\end{pfaufg}
% 2023-bebb-gk-B2.1j-m1 (Original 2023-bebb-gk-B2.1j, Schwere II, rang 28; herausgelöste Fassung: keine)
\begin{pfaufg}{GK~’23\\B2.1j}{5.}{}
Gegeben ist die in $\mathbb{R}$ definierte Funktion $f$ mit $f(x) = \tfrac{1}{3}(x^2 - 9) \cdot e^{\tfrac{x}{2}}$ und $f'(x) = \left(\tfrac{1}{6}x^2 + \tfrac{2}{3}x - \tfrac{3}{2}\right) \cdot e^{\tfrac{x}{2}}$. Die Abbildung zeigt die Draufsicht einer Tischplatte: Ihr Rand ist symmetrisch zu beiden Achsen, im III. Quadranten ist er der Graph $G$ von $f$, im II. Quadranten der Graph von $-f$. Berechne die Größe des Winkels, den diese beiden Randlinien im Punkt $(-3 \mid 0)$ einschließen.
\gzwei{\begin{abiksys}[xmin=-3.5,xmax=3.5,ymin=-3.5,ymax=3.5,le=0.5]\funktionab[-1.5]{(\x^2-9)*exp(\x/2)/3}{G}{-3}{0}\funktionab{-(\x^2-9)*exp(\x/2)/3}{}{-3}{0}\funktionab{(\x^2-9)*exp(-\x/2)/3}{}{0}{3}\funktionab{-(\x^2-9)*exp(-\x/2)/3}{}{0}{3}\end{abiksys}}{\karo{4}}
\fusshilfe{\gfuss{5}{48,1°}}
\end{pfaufg}

\pfblattende{10}{26 & 25 & 24 & 23 & 22 & 21 & 20 & 19 & 18 & 17}{ & B2.1d & B2.1i & B2.1j & B2.1e B2.2e B2.2f &  & B2.1f &  & B1.2b & B1.1c}%
{\pfausschnitt{Änderungsrate \pkt{} Tangente}{Winkel am Graphen}{Dreieck aus Tangente und Normale}}
\end{document}